\documentclass[11pt]{article}

\usepackage{amsmath, amssymb}
\usepackage{geometry}
\geometry{margin=1in}

\title{CSE400 — Fundamentals of Probability in Computing \\
Lecture L11--12: Transformation of Random Variables}
\author{Vansh Lilani (ID: AU2320146)}
\date{}

\begin{document}

\maketitle

\section{Topic Title}

\textbf{Transformation of Random Variables and Functions of Two Random Variables}

Topics covered:
\begin{itemize}
\item Transformation of random variables
\item Function of two random variables
\item Derived distribution example: $Z = X + Y$
\end{itemize}

\section{Definitions and Notation}

Let:
\begin{itemize}
\item $X$ be a continuous random variable with known PDF $f_X(x)$
\item $F_X(x)$ be the CDF of $X$
\end{itemize}

A transformation is defined as:
\[
Y = g(X)
\]

If invertible:
\[
x = g^{-1}(y)
\]

Goal:
\[
f_Y(y) \quad \text{and} \quad F_Y(y)
\]

For two random variables:
\[
Z = g(X,Y)
\]

Examples:
\[
Z = X + Y, \quad X - Y, \quad \frac{X}{Y}, \quad XY
\]

\section{Assumptions / Conditions}

\subsection{Monotonic Transformation}

The function $g(x)$ must be:
\begin{itemize}
\item monotonically increasing, or
\item monotonically decreasing
\end{itemize}

\section{Main Results / Theorems}

\subsection{Transformation Theorem (Single Random Variable)}

\subsubsection*{Step 1: CDF Transformation}

\[
F_Y(y) = P(Y \le y)
\]
\[
= P(g(X) \le y)
\]

If $g$ is increasing:
\[
F_Y(y) = P(X \le g^{-1}(y))
= F_X(g^{-1}(y))
\]

If $g$ is decreasing:
\[
F_Y(y) = P(X \ge g^{-1}(y))
= 1 - F_X(g^{-1}(y))
\]

\subsubsection*{Step 2: PDF Transformation}

\[
f_Y(y) = \frac{d}{dy}F_Y(y)
\]

\[
f_Y(y) = f_X(g^{-1}(y)) \left| \frac{dx}{dy} \right|
\]

where
\[
x = g^{-1}(y)
\]

\section{Proofs / Derivations}

\subsection{Function of Two Random Variables}

Let
\[
Z = X + Y
\]

\[
F_Z(z) = P(Z \le z)
\]
\[
= P(X + Y \le z)
\]

\subsubsection*{Integral Representation}

\[
F_Z(z) =
\int_{-\infty}^{\infty}
\int_{-\infty}^{z-y}
f_{XY}(x,y)\,dx\,dy
\]

\subsection{Leibniz Rule}

If
\[
G(x)=\int_{a(x)}^{b(x)} h(x,y)\,dy
\]

then

\[
\frac{d}{dx}G(x)
=
\frac{db(x)}{dx}h(x,b(x))
-
\frac{da(x)}{dx}h(x,a(x))
+
\int_{a(x)}^{b(x)}
\frac{\partial}{\partial x}h(x,y)\,dy
\]

\subsection{Differentiating the CDF}

\[
f_Z(z)=\frac{d}{dz}F_Z(z)
\]

\[
f_Z(z)=\int_{-\infty}^{\infty}
f_{XY}(z-y,y)\,dy
\]

Equivalent form:
\[
f_Z(z)=\int_{-\infty}^{\infty}
f_{XY}(x,z-x)\,dx
\]

\subsection{Independent Case (Convolution)}

If $X$ and $Y$ are independent:
\[
f_{XY}(x,y)=f_X(x)f_Y(y)
\]

\[
f_Z(z)=\int_{-\infty}^{\infty}
f_X(x)f_Y(z-x)\,dx
\]

This is called the \textbf{convolution integral}.

\section{Worked Examples}

\subsection{Example 1 — Transformation of RV}

Given:
\[
X \sim \text{Uniform}(-1,1)
\]

\[
f_X(x)=
\begin{cases}
\frac12, & -1<x<1 \\
0, & \text{otherwise}
\end{cases}
\]

Find PDF of:
\[
Y=\sin\left(\frac{\pi X}{2}\right)
\]

\subsubsection*{Solution}

\[
x=\frac{2}{\pi}\sin^{-1}(y)
\]

\[
\frac{dx}{dy}=\frac{2}{\pi}\frac{1}{\sqrt{1-y^2}}
\]

\[
f_Y(y)=f_X(x)\left|\frac{dx}{dy}\right|
\]

\[
f_Y(y)=\frac12\cdot\frac{2}{\pi}\frac{1}{\sqrt{1-y^2}}
\]

\[
f_Y(y)=\frac{1}{\pi\sqrt{1-y^2}}, \quad -1<y<1
\]

\subsection{Example 2 — Sum of Independent Normals}

\[
X \sim N(0,1), \quad Y \sim N(0,1)
\]

\[
f_Z(z)=\int_{-\infty}^{\infty}
\frac{1}{\sqrt{2\pi}}e^{-x^2/2}
\cdot
\frac{1}{\sqrt{2\pi}}e^{-(z-x)^2/2}
dx
\]

Result:
\[
Z \sim N(0,2)
\]

\subsection{Example 3 — Exponential RVs}

\[
X,Y \sim \text{Exponential}(\lambda)
\]

\[
f_X(x)=\lambda e^{-\lambda x}, \quad x>0
\]
\[
f_Y(y)=\lambda e^{-\lambda y}, \quad y>0
\]

\[
f_Z(z)=\int_0^z
\lambda e^{-\lambda x}
\lambda e^{-\lambda(z-x)}\,dx
\]

\[
=\lambda^2 e^{-\lambda z}\int_0^z dx
\]

\[
f_Z(z)=\lambda^2 z e^{-\lambda z}, \quad z>0
\]

\subsection{Example 4 — Difference of RVs}

\[
Z=X-Y
\]

\[
F_Z(z)=P(X-Y\le z)
\]

\[
=\int_{-\infty}^{\infty}
\int_{-\infty}^{z+y}
f_{XY}(x,y)\,dx\,dy
\]

\[
f_Z(z)=\frac{d}{dz}F_Z(z)
\]

\section{Additional Transformations Mentioned}

\[
Z=\frac{X}{Y}, \quad
R=\sqrt{X^2+Y^2}
\]

\end{document}